% TOP_PROBLEM: {"id": 7500117, "title": "Lebesgue Differentiation and Weak Weak König's Lemma", "category_id": 18, "category": {"id": 18, "name": "logic", "display_name": "Logic", "description": "Problems in mathematical logic, model theory, proof theory, and finite model theory.", "slug": "logic", "order_index": 18, "created_at": "2026-07-31T15:26:25.671Z"}, "status": "partially_solved"}
% FIRST_POSTED: null
% ENTRY_KIND: statement_only
% Imported from: batch20.tex; UnsolvedMath ID 7500117
\documentclass[11pt]{article}
\usepackage[a4paper,margin=25mm]{geometry}
\usepackage{amsmath,amssymb,mathtools}
\usepackage{fontspec,unicode-math}
\setmainfont{Latin Modern Roman}
\setsansfont{Latin Modern Sans}
\setmathfont{Latin Modern Math}
\usepackage{xurl,hyperref}
\hypersetup{colorlinks=true,linkcolor=blue,urlcolor=blue,pdftitle={UnsolvedMath Problem 7500117}}
\newcommand{\sref}[2]{\hyperlink{source:#1:#2}{[S#2]}}
\newcommand{\cataloguescope}[1]{\paragraph{Mathematical question.}#1}
\begin{document}
% UnsolvedMath ID 7500117; source catalogue code AMR-074-0117
\setcounter{section}{7500116}
\section{Lebesgue Differentiation and Weak Weak König's Lemma}\label{7500117}
{\small\sffamily UnsolvedMath ID 7500117 · Logic}
\subsection{Definitions and mathematical statement}

\paragraph{Shared notation.}$\mathbb N=\{1,2,\ldots\}$, and $\mathbb Z,\mathbb Q,\mathbb R,\mathbb C$ denote the integers, rationals, reals and complex numbers. Any other conventions are specified locally.


Work in second-order arithmetic over $\mathsf{RCA}_0$, the base theory with $\Delta^0_1$ comprehension and $\Sigma^0_1$ induction. The strength of a principle means which subsystems prove it and which subsystems it implies over that base; equivalence requires both implications. An $L^1$ function on $[0,1]^d$, $d\geq1$, is represented as a Cauchy limit of polynomials in the $L^1$ norm. The Lebesgue differentiation theorem says that for almost every $x$, $f(x)=\lim_{Q\ni x,\,\operatorname{diam}Q\to0}\mu(Q)^{-1}\int_Qf\,d\mu$, where $Q$ ranges over open cubes containing $x$ and $\mu$ is Lebesgue measure. The principle $\mathsf{WWKL}_0$ asserts that every binary tree $T\subseteq\{0,1\}^{<\omega}$ of positive measure has an infinite path; positive measure means $\inf_n |T\cap\{0,1\}^n|/2^n>0$.
\paragraph{Statement recorded in the supplied source.}
Does the Lebesgue differentiation theorem imply $\mathsf{WWKL}_0$ over $\mathsf{RCA}_0$? \sref{7500117}{2}
\subsection{Short English statement}
Does proving almost-everywhere differentiation for all integrable functions require the positive-measure binary-tree path principle?
\subsection{Sources}
\begin{description}
\item[\hypertarget{source:7500117:1}{S1}] ULAM.AI and the UnsolvedMath Contributors, UnsolvedMath: A Curated Collection of Open Mathematics Problems, record 7500117 (AMR-074-0117). CC BY 4.0. \url{https://huggingface.co/datasets/ulamai/UnsolvedMath}.
\item[\hypertarget{source:7500117:2}{S2}] Antonio Montalbán, Open Questions in Reverse Mathematics, author preprint dated 17 August 2010, Question 17, printed p. 7. \url{https://math.berkeley.edu/~antonio/papers/questionsRM.pdf}.
\end{description}
\label{end:7500117}
\end{document}
